\chapter{The Current Account}\label{ch:b3:current-account}

A country that spends more than it earns borrows from the rest of the world, and a
country that saves more than it invests lends. The \omterm{def:b3:current-account:accounts}{current account} records the
difference, year by year, and the international investment position records what
it adds up to. This chapter treats the \omterm{def:b3:current-account:accounts}{current account} as a decision --- how a
country spreads its consumption over time when it can borrow and lend abroad ---
and asks when the borrowing must stop, what happens when lenders stop it
suddenly, and why capital flows so little from rich countries to poor ones.

\begin{openingchart}{Current-account balances as a percent of world GDP, 1990--2024:
four economies, all other countries with data, and the sum over all countries
with data, which should be zero (130 countries in 1990, 189 in 2014, 169 in 2024).}{World Bank, World Development Indicators (BN.CAB.XOKA.CD,
NY.GDP.MKTP.CD); our computation}
\begin{tikzpicture}
\begin{axis}[omtime, width=0.72\linewidth, height=0.46\linewidth, xmin=1989, xmax=2025,
  ymin=-1.8, ymax=1.2, ylabel={Percent of world GDP}, legend columns=3,
  legend style={at={(0.5,1.03)}, anchor=south, font=\footnotesize}, unbounded coords=jump]
  \addplot[omDef, thick, mark=none] table[x=year, y=USA] {figdata/bachelor-3/12-current-account/out/ca_world.csv};
  \addplot[omThm, thick, mark=none] table[x=year, y=CHN] {figdata/bachelor-3/12-current-account/out/ca_world.csv};
  \addplot[black, thick, mark=none] table[x=year, y=DEU] {figdata/bachelor-3/12-current-account/out/ca_world.csv};
  \addplot[black, thick, dashed, mark=none] table[x=year, y=JPN] {figdata/bachelor-3/12-current-account/out/ca_world.csv};
  \addplot[omThm, thick, dashed, mark=none] table[x=year, y=rest] {figdata/bachelor-3/12-current-account/out/ca_world.csv};
  \addplot[gray, thick, dotted, mark=none] table[x=year, y=sum] {figdata/bachelor-3/12-current-account/out/ca_world.csv};
  \legend{United States, China, Germany, Japan (from 1996), all others, sum of all}
\end{axis}
\end{tikzpicture}
\end{openingchart}

% figdata: ca_world_points
The United States' deficit reached 1.5~\% of world GDP in 2006, and was 1.1~\% again
in 2024. China, Germany and Japan ran the matching surpluses, each a few tenths of
a percent of world output. And the sum over all countries, which accounting says
must be zero, was 0.36~\% of world GDP in 2019 and 0.35~\% in 2024: the world
appears to run a surplus with itself.

\begin{recall}
The grades 10--12 volume introduced the balance of payments; the Year 1 volume the
national accounts identity $Y = C + I + G + X - M$; the Year 2 volume intertemporal
choice and the Euler equation, $u'(c_1) = \beta(1 + r)u'(c_2)$.
\end{recall}

\begin{notation}[Symbols of this chapter]\label{not:b3:current-account:symbols}
$CA$ is the \omterm{def:b3:current-account:accounts}{current account}, $NX = X - M$ net exports, $S$ national saving, $I$
investment, $B$ (or $-D$) net foreign assets, and $r$ the world real interest rate.
Lower-case $nx$, $b$ and $d$ are ratios to GDP, and $g$ is the growth rate of GDP.
\end{notation}

\section{The balance of payments}\label{sec:b3:current-account:bop}

\begin{definition}[Current account; financial account]\label{def:b3:current-account:accounts}
\index{current account}\index{financial account}
The \emph{current account} records a country's net exports of goods and services,
its net primary income from abroad (interest, dividends, wages) and its net
secondary income (transfers). The \emph{financial account} records its net
acquisition of foreign assets minus its net incurrence of foreign liabilities.
With the small capital account (transfers of capital), the two sum to zero up to
errors and omissions: a current-account deficit is financed by a net sale of
assets to foreigners.
\end{definition}

\begin{definition}[Net international investment position; valuation effect]\label{def:b3:current-account:niip}
\index{net international investment position}\index{valuation effect}
The \emph{net international investment position} (NIIP) is the value of a
country's foreign assets minus its foreign liabilities at a date. A
\emph{valuation effect} is a change in it caused by changes in asset prices or
exchange rates rather than by transactions.
\end{definition}

\begin{proposition}[The current account as saving minus investment]\label{prop:b3:current-account:identity}
National saving minus investment equals the \omterm{def:b3:current-account:accounts}{current account}, and the \omterm{def:b3:current-account:accounts}{current
account}, plus the capital account and \omterm{def:b3:current-account:niip}{valuation effects}, equals the change in
the NIIP: $CA = S - I$ and $\Delta B = CA + KA + V$.
\end{proposition}

\begin{proof}
Gross national disposable income is $Y + NFI$, where $NFI$ is net primary and
secondary income from abroad. National saving is what is left after consumption:
$S = Y + NFI - C - G$. With $Y = C + I + G + NX$, $S - I = NX + NFI = CA$. Each
transaction in the \omterm{def:b3:current-account:accounts}{current account} has a counterpart in the \omterm{def:b3:current-account:accounts}{financial account},
which changes the stock of net foreign assets one for one; the capital account
adds its transfers, and price and exchange-rate changes $V$ revalue the stock.
\end{proof}

\begin{measure}[Why the world's current accounts do not add to zero]
Every export is someone's import, so the world's \omterm{def:b3:current-account:accounts}{current accounts} should sum to
zero. They do not, because each country measures its own side: shipments in transit
at the end of a year, income paid to holders who are hard to identify, services
that one side records and the other misses. In the chart the sum was a deficit of
0.6 to 1.0~\% of world GDP in the early 1990s, close to zero around 2005, and a
surplus of 0.15 to 0.6~\% since 2009. % figdata: ca_world
Part of that movement is coverage: the number of countries with data rises from
130 in 1990 to 189 in 2014, and a country missing from the sum takes its balance
with it --- Japan, absent before 1996, adds a surplus of 0.2~\% of world GDP when it
enters. Over the 133 countries with data in every year from 2000, the sum still
moves from a deficit of 0.6~\% in 2000 to a surplus of 0.5~\% in 2021.
% figdata: ca_balanced
The sign has changed with the same countries, so no single omission explains it; the discrepancy is a
reminder that each country's figure is an estimate, with an error of the size of
many countries' whole balances.
\end{measure}

\begin{measure}[Valuation effects in the investment position]
Because a country's assets and liabilities differ in currency and kind, their
values move apart even when no transaction occurs. A depreciation of the home
currency raises the home value of foreign-currency assets and leaves
home-currency liabilities unchanged. For the United States, which borrows mostly
in dollars and in safe short-term instruments and holds abroad equity and direct
investment, the return on its assets has exceeded the return it pays on its
liabilities --- a return difference within each class of assets, and a composition
effect of borrowing short and lending long. % ledger: b3ca:gr.privilege
\end{measure}

\section{The current account as a choice over time}\label{sec:b3:current-account:intertemporal}

\begin{definition}[Intertemporal approach to the current account]\label{def:b3:current-account:intertemporal}
\index{intertemporal approach to the current account}
The \emph{intertemporal approach to the \omterm{def:b3:current-account:accounts}{current account}} treats the \omterm{def:b3:current-account:accounts}{current account}
as the outcome of saving and investment decisions by forward-looking households
and firms that can borrow and lend abroad at the world interest rate.
\end{definition}

Take a small open economy that lives two periods, receives output $Y_1$ and $Y_2$,
and can borrow or lend at the world rate $r$. The household maximises $\ln C_1 + \beta
\ln C_2$ subject to $C_1 + C_2/(1 + r) = Y_1 + Y_2/(1 + r)$ (no initial assets, no
investment or government).

\begin{proposition}[Temporary and permanent shocks]\label{prop:b3:current-account:temporary}
If $\beta(1 + r) = 1$, consumption is the same in both periods and $CA_1 = Y_1 - \bar C$ with
$\bar C = (Y_1 + Y_2/(1 + r))/(1 + 1/(1 + r))$. A temporary fall in $Y_1$ alone produces a
current-account deficit of about half its size; a permanent fall of the same size
in both periods leaves the \omterm{def:b3:current-account:accounts}{current account} unchanged.
\end{proposition}

\begin{proof}
The Euler equation gives $1/C_1 = \beta(1 + r)/C_2$, so $C_1 = C_2 = \bar C$, and the budget
gives $\bar C$. If $Y_1$ falls by $\Delta$ with $Y_2$ fixed, $\bar C$ falls by $\Delta/(1 + 1/(1 + r))$,
slightly more than $\Delta/2$, so $CA_1 = Y_1 - \bar C$ falls by $\Delta(1 - (1 + r)/(2 + r)) = \Delta/(2 + r)$,
about half of $\Delta$. If both outputs fall by $\Delta$, $\bar C$ falls by $\Delta$, and $CA_1$ does not move.
\end{proof}

\begin{omfigure}
\begin{tikzpicture}[econ, x=0.55cm, y=0.55cm]
  \draw[econ axes] (0,0) -- (12,0) node[right] {$C_1$};
  \draw[econ axes] (0,0) -- (0,12) node[above] {$C_2$};
  \draw[curve, omDef] (0.5,11.571) -- (11.5,0.241) node[above right, font=\scriptsize] {slope $-(1 + r)$};
  \draw[curve, omThm, domain=3.5:11, samples=60] plot (\x, {37.39/(\x^1.03)});
  \fill (4.2,7.76) circle (2pt) node[left, font=\scriptsize] {endowment $(Y_1, Y_2)$};
  \fill (5.954,5.954) circle (2pt) node[above right, font=\scriptsize] {consumption};
  \draw[guide] (4.2,0) node[below, font=\scriptsize] {$Y_1$} -- (4.2,7.76);
  \draw[guide] (5.954,0) node[below, font=\scriptsize] {$C_1$} -- (5.954,5.954);
\end{tikzpicture}
\par\smallskip{\small The two-period small open economy after a temporary fall in
first-period output. The country borrows $C_1 - Y_1$ in period 1, a current-account
deficit, and repays it with interest in period 2.}
\end{omfigure}

The model predicts that a country runs deficits when output is temporarily low,
when its investment opportunities are good, and when it expects to be richer
later; and that permanent shocks pass straight to consumption. Its most direct
test is the link between saving and investment.

\section{Saving, investment and capital mobility}\label{sec:b3:current-account:fh}

\begin{definition}[Feldstein--Horioka puzzle]\label{def:b3:current-account:feldstein-horioka}
\index{Feldstein--Horioka puzzle}
The \emph{Feldstein--Horioka puzzle} is the finding that, across countries, average
investment rates move almost one for one with average saving rates, as if capital
did not cross borders, although financial markets are open.
\end{definition}

With perfect capital mobility, a country's investment depends on its returns, not
on its own saving, and the slope of investment on saving across countries should
be near zero. Feldstein and Horioka found a slope near one for rich countries; a
later estimate for 17 of them over 1974--1979 was 0.85: each extra dollar of
domestic saving raised domestic investment by about 85 cents. % ledger: b3ca:feld.retention

\begin{omchart}{Gross capital formation against gross national saving, country
averages 2000--2024 (percent of GDP), 122 countries with at least 20 years of data and
more than a million people, with the least-squares line.}{World Bank, World
Development Indicators (NY.GNS.ICTR.ZS, NE.GDI.TOTL.ZS); our computation}
\begin{tikzpicture}
\begin{axis}[omscatter, width=0.55\linewidth, height=0.45\linewidth, xmin=-10, xmax=60,
  ymin=0, ymax=50, xlabel={Saving (\% of GDP)}, ylabel={Investment (\% of GDP)}]
  \addplot table[x=saving, y=investment] {figdata/bachelor-3/12-current-account/out/fh.csv};
  \addplot[omThm, thick, no marks, sharp plot] table[x=saving, y=investment] {figdata/bachelor-3/12-current-account/out/fh_fit.csv};
\end{axis}
\end{tikzpicture}
\end{omchart}

% figdata: fh_stats
In a world sample of the last quarter-century, the slope is only 0.22 (standard
error 0.05), and saving explains 14~\% of the variation in investment. Capital is
more mobile, or the sample wider, than in the rich countries of the 1970s; the slope
also moves a great deal from decade to decade (\cref{exo:b3:current-account:7}).

\section{Sustainability}\label{sec:b3:current-account:sustainability}

\begin{definition}[External sustainability]\label{def:b3:current-account:sustainability}
\index{external sustainability}\index{no-Ponzi condition}
A path of net external debt is \emph{sustainable} if it satisfies the
\emph{no-Ponzi condition}: the present value of the debt far in the future is not
positive, $\lim_{T\to\infty} D_T/(1 + r)^T \le 0$. The country cannot forever pay the interest on its
debt by borrowing more.
\end{definition}

\begin{proposition}[The sustainable trade balance]\label{prop:b3:current-account:sustainable}
With net external debt $D$ at interest $r$, the \omterm{def:b3:current-account:sustainability}{no-Ponzi condition} holds if and only
if $D_0 \le \sum_{t \ge 1} NX_t/(1 + r)^t$. If GDP grows at $g < r$, the debt ratio $d$ stays
constant when net exports are $nx = d\,(r - g)/(1 + g)$ of GDP.
\end{proposition}

\begin{proof}
Each period $D_t = (1 + r)D_{t-1} - NX_t$. Dividing by $(1 + r)^t$ and summing from 1 to
$T$, $D_T/(1 + r)^T = D_0 - \sum_{t=1}^T NX_t/(1 + r)^t$; the limit is non-positive exactly when
$D_0$ does not exceed the present value of net exports. In ratios, $d_t = d_{t-1}(1 +
r)/(1 + g) - nx_t$, and $d_t = d_{t-1} = d$ requires $nx = d((1 + r)/(1 + g) - 1) = d(r - g)/(1 + g)$.
\end{proof}

A debtor whose economy grows faster than the interest rate can run trade deficits
for ever with a stable debt ratio; one that pays more than it grows must earn a
trade surplus. Sustainability is not solvency alone: lenders must also believe it.

\section{Sudden stops}\label{sec:b3:current-account:sudden}

\begin{definition}[Sudden stop]\label{def:b3:current-account:sudden-stop}
\index{sudden stop}
A \emph{sudden stop} is an abrupt fall in a country's net capital inflows, which
forces its current-account deficit to shrink by the same amount within a short
time.
\end{definition}

\begin{omchart}{Current-account balance before and after three sudden stops
(percent of GDP; year 0 is 1994 for Mexico, 1997 for Thailand, 2010 for
Greece).}{World Bank, World Development Indicators (BN.CAB.XOKA.GD.ZS)}
\begin{tikzpicture}
\begin{axis}[omtime, width=0.6\linewidth, height=0.4\linewidth, xmin=-3.3, xmax=3.3,
  ymin=-16, ymax=14, xtick={-3,...,3}, xlabel={Years from the stop},
  ylabel={Percent of GDP}, legend pos=north west, legend style={font=\footnotesize}]
  \addplot[omDef, thick, mark=*] table[x=k, y=MEX] {figdata/bachelor-3/12-current-account/out/sudden.csv};
  \addplot[omThm, thick, mark=square*] table[x=k, y=THA] {figdata/bachelor-3/12-current-account/out/sudden.csv};
  \addplot[black, thick, mark=triangle*] table[x=k, y=GRC] {figdata/bachelor-3/12-current-account/out/sudden.csv};
  \legend{Mexico, Thailand, Greece}
\end{axis}
\end{tikzpicture}
\end{omchart}

\begin{record}[Sudden stops of the 1990s]
% figdata: sudden
Mexico's \omterm{def:b3:current-account:accounts}{current account} was in deficit by 5.4~\% of GDP in 1994; in 1995, after the
peso collapsed, it was close to balance ($-0.4$~\%). Thailand's deficit of 8.0~\% of
GDP in 1996 became a surplus of 12.5~\% in 1998: a swing of more than 20 points in two
years. Neither country chose this adjustment. When lenders stop, the deficit must
close, through a fall in imports, which means a fall in spending and usually in
output, and a real depreciation that makes exports cheaper.
\end{record}

\begin{record}[The euro-area periphery, 2008--2013]
% figdata: sudden
Greece's deficit was 14.6~\% of GDP in 2008. Private inflows stopped in 2010, but
the adjustment was slower than Thailand's: official loans from other euro-area
governments and the European Central Bank's lending to Greek banks replaced
private capital. The deficit was still 10~\% of GDP in 2011 and fell to about 2.6~\% in
2012. A member of a currency union cannot depreciate its own currency; the real
depreciation had to come through lower wages and prices, and the fall in imports
through a deep recession.
\end{record}

\section{Why capital does not flow to poor countries}\label{sec:b3:current-account:lucas}

\begin{definition}[Lucas paradox]\label{def:b3:current-account:lucas-paradox}
\index{Lucas paradox}
The \emph{Lucas paradox} is the observation that, although the neoclassical model
predicts much higher returns to capital in poor countries, little capital flows
from rich countries to poor ones.
\end{definition}

\begin{proposition}[The Lucas calculation]\label{prop:b3:current-account:lucas}
If output per worker is $y = Ak^\alpha$ with the same $A$ in two countries, the ratio of
their marginal products of capital is $(y_{\mathrm{rich}}/y_{\mathrm{poor}})^{(1-\alpha)/\alpha}$.
\end{proposition}

\begin{proof}
$\mathrm{MPK} = \alpha A k^{\alpha - 1}$ and $k = (y/A)^{1/\alpha}$, so $\mathrm{MPK} = \alpha A^{1/\alpha}y^{(\alpha - 1)/\alpha}$.
With the same $A$ and $\alpha$, the ratio for the poor country to the rich one is
$(y_{\mathrm{poor}}/y_{\mathrm{rich}})^{(\alpha-1)/\alpha} = (y_{\mathrm{rich}}/y_{\mathrm{poor}})^{(1-\alpha)/\alpha}$.
\end{proof}

% figdata: lucas_stats
With the Penn World Table, output per worker in the United States was 7.18 times
India's in 2019; with $\alpha = 0.4$, the marginal product of capital in India would be
$7.18^{1.5} = 19.2$ times the American one. Lucas's own comparison, for 1988, gave a
factor of about 58. % ledger: b3ca:akv.lucas
Returns that differ by such factors would pull all new investment to poor
countries. They do not.

\begin{omchart}{Net inflows of foreign direct investment (median over 2000--2024,
percent of GDP) against GDP per person (mean over 2000--2024, constant 2021 PPP dollars,
log scale), 151 countries of more than a million people.}{World Bank, World
Development Indicators (BX.KLT.DINV.WD.GD.ZS, NY.GDP.PCAP.PP.KD); our computation}
\begin{tikzpicture}
\begin{axis}[omscatter, xmode=log, log ticks with fixed point, width=0.6\linewidth,
  height=0.42\linewidth, xmin=500, xmax=150000, ymin=-2, ymax=32,
  xlabel={GDP per person}, ylabel={FDI inflows (\% of GDP)}]
  \addplot table[x=y, y=fdi] {figdata/bachelor-3/12-current-account/out/lucas.csv};
  \addplot[omThm, thick, mark=square*, mark size=2.5pt, sharp plot] table[x=y, y=fdi] {figdata/bachelor-3/12-current-account/out/lucas_bins.csv};
\end{axis}
\end{tikzpicture}
\end{omchart}

% figdata: lucas_bins
The squares are the medians of the five quintiles of income: poor countries
received FDI of about 2.3~\% of their GDP, the richest fifth 2.7~\%, the middle about
3.3~\%. Measured against their size, poor countries do not draw more capital. The
explanations are those of growth accounting: technology and human capital differ
($A$ is not the same), and investors fear for their property. During 1970--2000, low
institutional quality was the leading explanation of the missing flows.
% ledger: b3ca:akv.lucas
Among developing countries, net capital flows even go more to countries that
invest and grow less, the opposite of what faster productivity growth should
attract. % ledger: b3ca:gj.allocation

\begin{model}{The intertemporal current account}
\assumes A small open economy with perfect capital mobility at a world interest
rate; households smooth consumption.
\predicts Deficits when output is temporarily low or investment opportunities are
good; no response to permanent shocks; investment unrelated to domestic saving;
capital flowing to where its return is highest.
\tested Saving and investment move together across countries, more so in the past
and among rich countries; capital flows little to poor countries, and among
developing countries goes to those that grow less. % ledger: b3ca:feld.retention, b3ca:gj.allocation
\breaks Borrowing limits and \omterm{def:b3:current-account:sudden-stop}{sudden stops}: lenders, not households, decide when the
deficit ends; returns to capital adjusted for risk and institutions are not
higher in poor countries.
\end{model}

\section{Exercises}\label{sec:b3:current-account:exercises}

\begin{exercise}[$\star$]\label{exo:b3:current-account:1}
In the opening chart, in which year was the United States' deficit largest
relative to world GDP? Which economies ran the matching surpluses in 2006, and how
large was the world discrepancy?
\end{exercise}

\begin{exercise}[$\star$]\label{exo:b3:current-account:2}
Record in Brenland's balance of payments: (a) a firm exports machines worth 100 and is
paid into its account at a Kestrian bank; (b) a household receives dividends of 10
from Kestrian shares; (c) a Kestrian fund buys Brenland government bonds worth 50;
(d) Brenland's government gives 5 in aid to a foreign country.
\end{exercise}

\begin{exercise}[$\star$]\label{exo:b3:current-account:3}
National saving is 22~\% of GDP and investment 25~\%. Find the \omterm{def:b3:current-account:accounts}{current account}. If the
budget deficit rises by 2~\% of GDP and private saving does not change, what happens
to it?
\end{exercise}

\begin{exercise}[$\star$]\label{exo:b3:current-account:4}
What pattern in a country's \omterm{def:b3:current-account:accounts}{current account} and output would contradict
consumption smoothing?
\end{exercise}

\begin{exercise}[$\star\star$]\label{exo:b3:current-account:5}
With $\ln C_1 + \beta\ln C_2$, $r = 5~\%$, $\beta = 1/1.05$, $Y_2 = 100$ and $Y_1 = 90$ (a temporary fall),
find $C_1$, $C_2$ and the \omterm{def:b3:current-account:accounts}{current account} in both periods. Do the same if $Y_1 = Y_2 = 90$.
\end{exercise}

\begin{exercise}[$\star\star$]\label{exo:b3:current-account:6}
A country's net external debt is 60~\% of GDP, the interest rate 4~\% and growth 3~\%.
What trade balance keeps the debt ratio constant? And if growth falls to 1~\%?
\end{exercise}

\begin{exercise}[$\star\star$]\label{exo:b3:current-account:7}
Estimate the Feldstein--Horioka slope on WDI data for 1990--1999, 2000--2011 and
2012--2024, using countries with every year of data and more than a million people.
Is the coefficient stable?
\end{exercise}

\begin{exercise}[$\star\star$]\label{exo:b3:current-account:8}
A country's \omterm{def:b3:ricardian-trade:terms-of-trade}{terms of trade} fall permanently by 10~\%. What does the intertemporal
model predict for its \omterm{def:b3:current-account:accounts}{current account}? And if markets expect the fall to reverse in
two years?
\end{exercise}

\begin{exercise}[$\star\star$]\label{exo:b3:current-account:9}
What observations would show that capital is perfectly mobile between two
countries? Name two.
\end{exercise}

\begin{exercise}[$\star\star\star$]\label{exo:b3:current-account:10}
Output per worker in the United States is 7.18 times India's. Find the ratio of
marginal products of capital with $\alpha = 0.4$ and with $\alpha = 1/3$. If workers in the
United States have 2.5 times the human capital of Indian workers and $y = Ak^\alpha h^{1-\alpha}$,
what is the ratio?
\end{exercise}

\begin{exercise}[$\star\star\star$]\label{exo:b3:current-account:11}
A country has a current-account deficit of 6~\% of GDP when inflows stop. Net exports
must rise by the full 6 points. If each 1~\% real depreciation raises net exports by
0.25~\% of GDP, how large a real depreciation is needed? If instead imports, at 30~\%
of GDP, must fall with spending and the import share of spending is 0.3, by how much
must spending fall?
\end{exercise}

\begin{exercise}[$\star\star\star$]\label{exo:b3:current-account:12}
Find current-account reversals in the WDI data since 1990: a deficit of at least
4~\% of GDP that shrinks by at least 4 points within two years, in countries of more
than a million people. How many are there, and what happens to GDP growth around
them?
\end{exercise}

\section{Problem: Brenland's Borrowing Years}\label{sec:b3:current-account:problem}

\begin{problem}[{Weekend problem --- Brenland's borrowing years}]\label{pb:b3:current-account:borrowing}
Brenland's net exports and GDP come from its national accounts. Assume that its net
external debt was 45{,}000 million dalens at the end of 1999, that it pays interest
on it at the policy rate plus one point, and that there are no transfers and no
valuation changes. Then $CA_t = NX_t - i_tD_{t-1}$ and $D_t = D_{t-1} - CA_t$:
% figdata: borrow_points
\begin{center}\small
\pgfplotstabletypeset[col sep=comma, columns={year,gdp,nx,i,interest,ca,debt_pct},
  columns/year/.style={column name={Year}, int detect, column type=l},
  columns/gdp/.style={column name={GDP}, fixed, precision=0, 1000 sep={\,}},
  columns/nx/.style={column name={$NX$}, fixed, precision=0, 1000 sep={\,}},
  columns/i/.style={column name={$i$ (\%)}, fixed, fixed zerofill, precision=2},
  columns/interest/.style={column name={Interest}, fixed, precision=0, 1000 sep={\,}},
  columns/ca/.style={column name={$CA$}, fixed, precision=0, 1000 sep={\,}},
  columns/debt_pct/.style={column name={$D$ (\% of GDP)}, fixed, fixed zerofill, precision=2},
  every head row/.style={after row=\hline}, every last row/.style={after row=\hline}]
  {figdata/bachelor-3/12-current-account/out/borrow_points.csv}
\end{center}
(Millions of dalens.) Over 2021--2025 the interest rate averaged 5.90~\%, and nominal
GDP grew by 4.76~\% a year on average over 2016--2025. % figdata: borrow_stats

\medskip\noindent\textbf{Part I --- the balance of payments.}
\begin{enumerate}
\item Compute the 2000 \omterm{def:b3:current-account:accounts}{current account} from the table's other columns.
\item In which listed years did Brenland run a trade surplus but a current-account
deficit? Why?
\item Express the 2025 \omterm{def:b3:current-account:accounts}{current account} as a percent of GDP. How was it financed?
\item Check the change in debt from 2009 to 2010.
\item The debt ratio rose from 30.3~\% in 2008 to 38.6~\% in 2020. Through which two
channels can a debt ratio rise?
\item What would a 10~\% depreciation of the dalen do to the NIIP if half of the debt
were in kests? (Give the effect in percent of 2025 GDP.)
\end{enumerate}

\medskip\noindent\textbf{Part II --- the 2009 export collapse.}
\begin{enumerate}[resume]
\item Real GDP fell by 3.3~\% in 2009. In a two-period model with $\ln C_1 + \beta\ln C_2$,
$\beta(1 + r) = 1$, $r = 3~\%$, normal output 100 and $Y_1 = 96.7$, find consumption.
\item Find the \omterm{def:b3:current-account:accounts}{current account} in period 1, in percent of normal output.
\item Compare it with the change in Brenland's \omterm{def:b3:current-account:accounts}{current account} from 2008 to 2009.
\item What would the model predict had the fall been permanent?
\item Which assumption of the model fits Brenland in 2009 poorly?
\item In 2020 real GDP fell by 4.85~\%. Was the change in the deficit from 2019 to
2020 consistent with a temporary shock?
\end{enumerate}

\medskip\noindent\textbf{Part III --- sustainability.}
\begin{enumerate}[resume]
\item Write the law of motion of the debt ratio.
\item Find the trade balance, in percent of GDP, that keeps the 2025 debt ratio
constant with the average interest rate and growth rate.
\item Compare it with the 2025 trade balance (0.28~\% of GDP). Is the debt ratio
rising?
\item If the interest rate fell to 4~\%, what trade balance would stabilise the
ratio?
\item Does the \omterm{def:b3:current-account:sustainability}{no-Ponzi condition} require Brenland to repay its debt?
\item Why might lenders stop before the condition is violated?
\end{enumerate}

\medskip\noindent\textbf{Part IV --- a sudden stop.}
\begin{enumerate}[resume]
\item If lenders refuse any new lending in 2026, what \omterm{def:b3:current-account:accounts}{current account} must Brenland
run?
\item With 2025's interest bill, what trade balance does that require, in percent of
GDP? How large is the swing from 2025?
\item If each 1~\% real depreciation of the dalen raises net exports by 0.3~\% of GDP,
how large a depreciation is needed?
\item State the result: the trade surplus, in percent of GDP, that stabilises
Brenland's external debt.
\end{enumerate}
\end{problem}
